\documentclass[11pt]{article}
\usepackage[final]{acl}
\usepackage{times}
\usepackage{latexsym}
\usepackage[T1]{fontenc}
\usepackage[utf8]{inputenc}
\usepackage{microtype}

\title{Medical Question Answering on PubMedQA:\\Analysing Text Representations and Classifier Design}

\author{Group 16 \\
  EMATM0067 -- Introduction to AI and Text Analytics \\
  University of Bristol}

\begin{document}
\maketitle

\begin{abstract}
Biomedical question answering requires classifying whether published evidence supports a yes, no, or maybe answer to a clinical research question. We investigate this three-class classification problem on the PubMedQA dataset, systematically exploring two design axes: (1)~text representation, comparing TF-IDF n-grams, BioBERT sentence embeddings, and varying input compositions; and (2)~classifier design, comparing Logistic Regression, calibrated Linear SVM, Multinomial Naive Bayes, SMOTE-augmented models, a soft voting ensemble, zero-shot LLM classification, and fine-tuned BioBERT. Our best model, a fine-tuned BioBERT classifier with class-weighted cross-entropy loss, achieves a test macro-F1 of 0.4108. Extensive error analysis reveals that the \textit{maybe} class is misclassified 81.8\% of the time, and 24\% of test samples are misclassified by all models, representing genuinely ambiguous cases. A quality predictor analysis confirms that surface-level text features cannot reliably predict answer correctness.
\end{abstract}

\end{document}